
\documentclass[11pt]{article}

\usepackage{amsfonts}
\usepackage{amsmath}
\usepackage{amscd}
\usepackage{amssymb}
\usepackage[tableposition=top]{caption}
\usepackage{ifthen}
\usepackage{indentfirst}
\usepackage{natbib}
\usepackage{rotating}
\usepackage[colorlinks=true,allcolors=blue]{hyperref}
\usepackage{url}
\usepackage{doi}
\usepackage{tikz-cd}

\DeclareMathOperator{\idfun}{id}
\DeclareMathOperator{\isContr}{isContr}
\DeclareMathOperator{\isProp}{isProp}
\DeclareMathOperator{\isSet}{isSet}
\DeclareMathOperator{\dom}{dom}
\DeclareMathOperator{\cod}{cod}
\DeclareMathOperator{\qinv}{qinv}
\DeclareMathOperator{\binv}{biinv}
\DeclareMathOperator{\isEquiv}{isequiv}
\DeclareMathOperator{\refl}{refl}
\DeclareMathOperator{\inl}{inl}
\DeclareMathOperator{\inr}{inr}
\DeclareMathOperator{\successor}{succ}
\DeclareMathOperator{\inductor}{ind}
\DeclareMathOperator{\recursor}{rec}

\newcommand{\sloop}{\text{loop}}
\newcommand{\sbase}{\text{base}}
\newcommand{\ssurf}{\text{surf}}

\newcommand{\set}[1]{\{\,#1\,\}}
\newcommand{\proptrunc}[1]{\lVert #1 \rVert}
\newcommand{\termtrunc}[1]{\lvert #1 \rvert}
\newcommand{\bigproptrunc}[1]{\left\lVert #1 \right\rVert}

\newcommand{\ints}{\mathbb{Z}}
\newcommand{\nats}{\mathbb{N}}
\newcommand{\real}{\mathbb{R}}
\newcommand{\universe}{\mathcal{U}}
\newcommand{\rats}{\mathbb{Q}}

\newcommand{\fatdot}{\,\cdot\,}

\let\emptyset=\varnothing

\newcommand{\REVISED}{\begin{center} \LARGE REVISED DOWN TO HERE \end{center}}
\newcommand{\MOVED}[1][equation]{\begin{center} [#1 moved] \end{center}}

\begin{document}

\title{Programming in Homotopy Type Theory}
\author{Charles J. Geyer}
\maketitle

\tableofcontents

\section{License and Source Code}

This work is licensed under a Creative Commons
Attribution-ShareAlike 4.0 International License
(\url{http://creativecommons.org/licenses/by-sa/4.0/}).

The \LaTeX\ source for this document is
at \url{https://github.com/cjgeyer/Survey}.

\pagebreak[3]
\section{Introduction}

Homotopy type theory (HoTT) can be thought of as explained in
\citet[hereinafter referred to as the HoTT book]{hott-book},
\defcitealias{hott-book}{the HoTT book}
a language for formalizing mathematics and hence a foundation for mathematics.
But it can also be thought of as a functional programming language.

A book presenting this point of view for Martin-L\"{o}f type theory (MLTT,
\citealp{martin-lof-mltt}) is \citet*{programming-mltt} but that book does
not seem to use exactly the same version of MLTT that HoTT uses.  Hence this
document.

In particular, the notion of type does not occur in \citet{programming-mltt}
until Chapter~19 (137 pages in).  Before that it uses sets.  This is
(apparently) very unlike HoTT.  So we will have to translate
\citet{programming-mltt} into HoTT.

We start with \citep[p.~1]{programming-mltt}
\begin{quotation}
In recent years several formalisms for program construction have been
introduced. One such formalism is the type theory developed by
Per Martin-L\"{o}f. It is well suited as a theory for program construction
since it is possible to express
both specifications and programs within the same formalism. Furthermore, the
proof rules can be used to derive a correct program from a specification as
well as to verify that a given program has a certain property.  This book
contains an
introduction to type theory as a theory for program construction.

As a programming language, type theory is similar to typed functional
languages such as Hope and ML [citations omitted], but a major difference
is that the evaluation of a well-typed program always terminates.
In type theory it is also possible
to write specifications of programming tasks as well as to develop provably
correct programs. Type theory is therefore more than a programming language and
it should not be compared with programming languages, but with formalized
programming logics such as LCF and PL/CV [citations omitted].
\end{quotation}

Section~1.1 of \citetalias{hott-book} gives multiple ways of talking about
the fundamental judgement $a : A$ read neutrally as $a$ is a term of type $A$.
Chapter~2 of \citet{programming-mltt} gives some others.
These readings are summarized in Table~\ref{tab:ways}.
\begin{table}
\caption{Interpretations of $a : A$.  Alternative for term is token.
  Alternatives for proof are witness, evidence, or certificate.}
\label{tab:ways}
\begin{center}
\begin{tabular}{ll}
term & type \\
\hline
proof & proposition \\
element & set \\
point & space \\
program & specification \\
solution & problem
\end{tabular}
\end{center}
\end{table}

When we think of HoTT as a language for formalizing mathematics, we are
primarily interested in proofs.  In this reading, if $A$ is a type considered
as a proposition, then proving that proposition is to show that $A$ is
inhabited (there exists $a : A$).
When we think of HoTT as a computer programming language, we are interested
in running programs.  So this means computing (or constructing) a particular
$a : A$, not just showing that such exists.

Of course, one way to demonstrate $a : A$ is to actually construct $a$.
So these are intertwined.  You cannot talk about one without the other.

One important point of the quotation above is that HoTT (or MLTT) is not
a Turing-complete computer language.  There is no halting problem.
All well-typed programs halt.  Moreover, all well-typed programs are
proved correct.  If it type checks, then it is correct.  Correctness is
just $a : A$.   The program satisfies its specification.

So one can think there are two languages: the computer programming language
and the program specification language, but these are the same language, HoTT
or MLTT, thought of in two different ways.

As a programming language, HoTT or MLTT is a pure functional language.
It does not have assignment or other side effects.  It does have functions
as first-class objects, and it does have higher-order functions, those that
have other functions as arguments or values.  Thus it contains higher-order
logic.

\section{Fundamentals}

\subsection{Function Types}

For any types $A$ and $B$, there is a function type $A \to B$.
If this type is inhabited, the inhabitant $f : A \to B$ is a function
and a program.  Given $x : A$ we can run the program to obtain $f(x) : B$.

Note that this is unlike conventional Turing-complete computer languages.
In MLTT or HoTT we cannot write a function that fails to terminate (does
an infinite loop) or fails with an error in any other way.

Note also that this is unlike functions in classical mathematics (using the
law of excluded middle (LEM) or the axiom of choice (AC)).
HoTT functions are given
by a formula $f(x) = \Phi$, where $\Phi$ is a HoTT formula that may contain
$x$ (if it does not contain $x$, then this is a constant function).
In classical mathematics, LEM and AC assert the existence
of functions for which there is no rule, hence no way the function can
be evaluated.

In order for a function definition to be valid, it must type-check.
The user or a proof assistant much check that $f(x) : B$ whenever $x : A$.

Evaluation of the function is just done by substitution.  To evaluate $f(y)$
for some $y : A$, substitute $y$ for $x$ in $\Phi$ taking the usual care that
$\Phi$ cannot contain $y$ as a bound variable.  (If $\Phi$ did contain $y$ as
a bound variable, then such occurrences must be renamed before the
substitution is done; \citetalias[Section~1.2]{hott-book}.  Of course,
this is getting us a bit ahead of ourselves because we have not yet discussed
the quantifiers $\prod$ and $\sum$ that bind variables.)

A function may be anonymous, having no name (like $f$).  This is called
$\lambda$-abstraction in computer science, and such functions are sometimes
called lambdas rather than anonymous functions.  Mathematicians generally
write such as $x \mapsto \Phi$.
The computer science notation is $\lambda x. \Phi$.
The $\lambda$ operator also binds variables (like $\prod$ and $\sum$ mentioned
in the preceding paragraph).

So now we have a real computer language that can calculate anything
any other computer language can.  This is because functions are enough
to express anything that can be computed (the Church-Turing thesis
and Church's lambda calculus).  So in a sense, we do not need the rest
of HoTT.  The other type formation rules are just for convenience.
It is known that they could be defined in terms of functions only.
(Appendix~A.1 of \citet{hott-book} does this.)

One thing that should be pointed out is currying, how to make a function
of multiple arguments.  All functions $A \to B$ have only one argument
of type $A$.  But if $B$ is itself a function type, say $B \equiv C \to D$,
then $f(x)$ is a function $C \to D$, which can be evaluated at $y : C$
giving $f(x)(y) : D$.  And so on for even more arguments.
The type of $f$ here is $A \to (C \to D)$.
To simplify this further, $\to$ is understood as right-associative,
so this can be written $A \to C \to D$.
(This is the convention in functional programming languages
in computer science and also in HoTT.)

Currying is how we build functions of many arguments out of functions of
single arguments.

Usually no notational fuss is made about currying.
If $f : A \to B \to C \to D$, then we should write $f(a)(b)(c)$ where
$a : A$, $b : B$, and $c : C$ to obtain $f(a)(b)(c) : D$.  But we often
consider this a function of three arguments and write $f(a, b, c)$ instead,
like in conventional mathematics.

\subsection{Dependent Function Types}

For any type $A$, if $B(x)$ is a type for all $x : A$, then there is
the dependent function type
$$
   \prod_{x : A} B(x)
$$
If $f$ is a term of this type, then $f$ is a function given by a formula
$f(x) = \Phi$, and in order to type check, either we or a proof assistant
must show that $f(x) : B(x)$ for $x : A$.
If we have such a function, then
given $x : A$ we can run the program to obtain $f(x) : B(x)$.
So dependent function types are just like non-dependent function types,
except that the codomain depends on the argument.

Currying works the same for dependent function types as for ordinary function
types.

In case the dependent type family $B(x) : x : A$ is actually constant,
that is $B(x)$ does not depend on $x$, then the dependent function type is
the same as the ordinary function type, that is,
$$
   \prod_{x : A} B \equiv A \to B
$$
where $\equiv$ means definitionally equal \citepalias[Section~1.1]{hott-book}.

The expression $\prod_{x : A} \Phi$ binds free occurrences of $x$ in $\Phi$.
Thus we cannot substitute $x$ in this expression without renaming the $x$
that is already bound in it.  First we rename $x$, and then we substitute.

\subsection{The Empty Type}

The empty type \citepalias[Section~1.7]{hott-book} is the type $0$ having
no terms.  This type has a formation rule: 0 is a type.
It has no constructors (rules that make terms).  Thus, in a sense there is
nothing to do.  There are no programs that satisfy the specification $0$.

But we are not finished with the rules for this type: it has eliminators
(rules for using terms).
The recursor for the empty type says for any type $A$ there is an
inhabited function type $0 \to A$.  Another way to say this is
$$
   \recursor_0 : \prod_{C : \universe} (0 \to C)
$$
Now this does give us something to compute.  The specification $0 \to C$
always has one program, the \emph{empty function}, which satisfies it.
This function has a vacuous rule, if $f : 0 \to C$, then we can evaluate
$f(x)$ for all $x : 0$ but there are no such $x$.

So we have a function $f : 0 \to C$, which is a program.  But this program
cannot do anything, we cannot evaluate $f(x)$ because there are no valid
$x$ that can serve as arguments to this function.

The induction principle for this type says the same thing for
dependent type families
$$
   \inductor_0 : \prod_{C : 0 \to \universe} \prod_{x : 0} C(x)
$$
for any dependent type family $C : 0 \to \universe$ there is
a dependent function type $\prod_{x : 0} C(x)$, which again is the
empty function having the vacuous rule and no valid arguments.

\subsection{The Unit Type}

The unit type \citepalias[Sections~1.5 and~2.8]{hott-book} is the type 1
having one term, denoted $\star$.  This type has a formation rule: 1 is a type.
It has one constructor
$$
   \star : 1
$$
There is a program $\star$ that satisfies the specification $1$.
But it doesn't compute anything.  It is just a defined constant
(the program does not define a constant something-or-other; it is a constant).

The recursor for the unit type says for any inhabited type $C$ there is an
inhabited function type $1 \to C$.  Another way to say this is
$$
   \recursor_1 : \prod_{C : \universe} (C \to 1 \to C)
$$
having defining equation
\begin{equation} \label{eq:recursor-1-definition}
   \recursor_1(C, c, \star) = c
\end{equation}
Decoding this (and getting a bit ahead of ourselves)
\begin{itemize}
\item $\recursor_1$ is a dependent product type, hence a function whose
    argument is the subscript of the $\prod$, that is,
\item $\recursor_1(C)$
    is a term of type $C \to 1 \to C$.  And $\to$ is right-associative
    by default so this means $C \to (1 \to C)$.  Hence
\item $\recursor_1(C)(c)$ is for $c : C$ a function $1 \to C$.  Hence
\item $\recursor_1(C)(c)(\star)$ is a term of type $C$ given by
    \eqref{eq:recursor-1-definition}.
\item Except for the case $C \equiv 0$ in which case
    $\recursor_1(C)(c)$ does not exist because there is no $c : 0$.
    In this case $\recursor_1(0) : 0 \to 1 \to 0$ is inhabited because
    $1 \to 0$ is equivalent to $0$ (has no terms), so $\recursor_1(0)$ is
    the empty function.
\end{itemize}
In short, for any type $C$ and any $c : C$
there exists $f : 1 \to C$ defined by $f(\star) = c$.
The equality of $\recursor_1(C, c, \star)$ and $\recursor_1(C)(c)(\star)$
is currying.

So we have a function $f : 1 \to C$, which is a program.  And this program
can do something, we can evaluate $f(\star)$ obtaining $c$.

The induction principle for this type says the same thing for
dependent type families
$$
   \inductor_1 : \prod_{C : 1 \to \universe} C(\star) \to \prod_{x : 1} C(x)
$$
with defining equation
\begin{equation} \label{eq:inductor-1-definition}
   \inductor_1(C, c, \star) = c
\end{equation}
So this says for any dependent type family $C(x)$, $x : 1$
\begin{itemize}
\item $\inductor_1(C)$ is a function $C(\star) \to \prod_{x : 1} C(x)$.
    Hence for $c : C(\star)$
\item $\inductor_1(C)(c)$ is a term of type $\prod_{x : 1} C(x)$.  Hence
\item $\inductor_1(C)(c)(x)$ is a term of type $C(x)$ given
    by \eqref{eq:inductor-1-definition}.
\end{itemize}
But now this does not make obvious sense, we need a theorem given in
Section~1.5 of \citetalias{hott-book} that  uses this induction principle.
Read it as a proof of a theorem rather than a program.
If we take the case $C(x) :\equiv (x = \star)$, then the induction principle
says $c : C(\star)$ implies
$$
   \prod_{x : 1} (x = \star)
$$
So this type has exactly one term $\star$.
Or, if we take the homotopy reading,
this space has exactly one path component.
Thus we have proved (following \citetalias{hott-book}) that this type has only
one term (up to propositional equality).

Now the principle of indiscernibility of identicals
\citepalias[Section~1.12]{hott-book} says that $x = \star$ implies that
$C(x)$ and $C(\star)$ are equivalent types, hence identical types by the
univalence axiom.

So now if we have the program specification $\prod_{x : 1} C(x)$ and
$c : C(\star)$, there is a function $f$, which is a program satisfying this
type satisfying $f(\star) = c$ and also $f(x) = c$ for any $x : 1$.
And we can execute this function by evaluating $f(x)$ to obtain $c$.

\subsection{Product Types}

For any types $A$ and $B$, there is a product type $A \times B$.
If this type is inhabited, the inhabitant $p : A \times B$ is an ordered
pair $(a, b)$ with $a : A$ and $b : B$.

The product type has one constructor.  Given $a : A$ and $b : B$, the
constructor gives $(a, b) : A \times B$.

The eliminator (recursor) of the product type says that a function
$f : A \times B \to C$ is defined by currying is definitionally equivalent
to a function $g : A \to B \to C$ with the equivalence
being $f((x, y)) = g(x)(y)$.

The dependent eliminator (induction principle) works the same way for
dependent function types.

As special cases we can derive the projection functions that extract the
terms that went into the pair.
\begin{align*}
   \text{fst} : A \times B \to A
   \\
   \text{snd} : A \times B \to B
\end{align*}
having the rules
\begin{align*}
   \text{fst} : (a, b) \mapsto a
   \\
   \text{snd} : (a, b) \mapsto b
\end{align*}

\begin{thebibliography}{}

\bibitem[Martin-L\"{o}f(1975)]{martin-lof-mltt}
Martin-L\"{o}f, P. (1975).
\newblock An intuitionistic theory of types: predicative part.
\newblock In: H. E. Rose, J. C. Shepherdson (eds.),
    \emph{Logic Colloquium ‘73, Proceedings of the Logic Colloquium,
    Studies in Logic and the Foundations of Mathematics}, \textbf{80}, 73--118.
\newblock Elsevier.
\newblock \doi{10.1016/S0049-237X(08)71945-1}.

\bibitem[Nordstr\"{o}m, et al.(1990)Nordstr\"{o}m, Petersson,
    and Smith]{programming-mltt}
Nordstr\"{o}m, B., Petersson, K., and Smith, J.~M. (1990)
\newblock Programming in Martin-L\"{o}f's Type Theory: An Introduction.
\newblock Clarendon Press, New York.

\bibitem[Rijke(2025)]{rijke}
Rijke, E. (2025).
\newblock \emph{Introduction to Homotopy Type Theory}.
\newblock Cambridge University Press.
\newblock \doi{10.1017/9781108933568}.
\newblock A preprint is available at
    \href{https://arxiv.org/abs/2212.11082}{Ar$\chi$iv}.

\bibitem[The Univalent Foundations Program(2013)]{hott-book}
The Univalent Foundations Program (2013).
\newblock \emph{Homotopy Type Theory: Univalent Foundations of Mathematics}.
\newblock Institute for Advanced Study, Princeton, NJ.
\newblock \url{https://homotopytypetheory.org/book}.

\end{thebibliography}

\end{document}

